\subsection{The cubic action and the physical curvature}
\label{subsec:representative-cubic}

The compact cubic block illustrates the curvature mechanism behind the
three-zero bound. At the scale \eqref{deg:cubic-chart}, its axis
polynomial is
\[
 P(\xi)=1-U\xi+3G\xi^2+\xi^3,
 \qquad U\ge0,
\]
with \(U,G\) in a fixed compact set. On a strict through interval,
let \(u>0\) and \(-r<0\) be the incoming and outgoing axis points.
The limiting action equation is
\begin{equation}\label{intro:compact-action}
 \int_0^u\frac{\xi}{P(\xi)}\,d\xi
 =\int_0^r\frac{\xi}{P(-\xi)}\,d\xi,
 \qquad b=-u^2/2,\qquad S(b)=-r^2/2.
\end{equation}
On the traversed segment \(P>0\). Differentiating gives
\(r_u=uP(-r)/(rP(u))>0\) and
\begin{equation}\label{intro:compact-action-curvature}
 S'=\frac{P(-r)}{P(u)},\qquad
 \frac{P(u)}{u+r}\frac{S''}{S'}=3-\frac{U}{ur},\qquad
 (ur)_b=-\frac ru-\frac{uP(-r)}{rP(u)}<0.
\end{equation}
These are the cubic specialization of
\eqref{neg:action-curvature}--\eqref{neg:entry-exit-monotonicity}.
The term containing \(G\) cancels. At a curvature zero, \(U=3ur\),
so the derivative of the middle expression with respect to \(b\)
is strictly negative. Thus every curvature zero is a downward crossing,
and there can be at most one on a connected interval.

The actual input is an affine raw section coordinate. To compare its
curvature with the action formula, write the two separate outward maps as
\[
 \Gamma_\pm=a_\pm+c_\pm\phi_\pm,\qquad c_\pm>0,
 \qquad \widehat\Pi=\phi_-\circ S_{\rm act}\circ\phi_+^{-1}.
\]
Here \(S_{\rm act}\) is the actual matched transition in the two normal
section labels. The centers and gains are retained exactly. In these
coordinates the second-derivative chain rule is
\begin{equation}\label{intro:completed-curvature-chain}
 \phi_+'\frac{\widehat\Pi''}{\widehat\Pi'}
 =\frac{S_{\rm act}''}{S_{\rm act}'}
  +S_{\rm act}'\frac{\phi_-''}{\phi_-'}
  -\frac{\phi_+''}{\phi_+'}.
\end{equation}
Each section derivative uses its own quadratic label. The actual gains
do not enter this identity. Returning to the physical affine input
multiplies its second and third derivatives by
\(c_-/c_+^2\) and \(c_-/c_+^3\), so their signs are preserved.

On its bounded normal-section interval,
Proposition~\ref{deg:prop-compact-cubic} proves the weighted identity
\begin{equation}\label{intro:actual-cubic-curvature}
 \phi_+'(b)\frac{P(u)}{u+r}
              \frac{\widehat\Pi''}{\widehat\Pi'}
 =3-\frac{U}{ur}+E(u,r),\qquad
 \|E\|_{C^1_{u,r}}\le C\tau_{\rm cub}(1+|\log\tau_{\rm cub}|).
\end{equation}
It includes both section terms in
\eqref{intro:completed-curvature-chain}. Near a possible zero the
principal derivative with respect to \(u\) has a margin
\(c(1+r_u)\), whereas the error derivative along the matched graph is
at most \(C\tau_{\rm cub}(1+|\log\tau_{\rm cub}|)(1+r_u)\).
The physical input \(\phi_+(-u^2/2)\) decreases with \(u\).
Consequently the actual curvature retains the same downward crossing.
The bounded-action comparison overlaps this description on the same
map. The rest of the admitted interval has bounded physical first
derivative; choosing \(m\) last makes \(F'>1/2\) there.
Lemma~\ref{sharp:upper-curvature-count} then gives three zeros on the
whole interval.

Two estimates make this calculation compatible with singular
coefficient limits. First, the exterior continuations carry the common
first-variation factor \(\tau_{\rm cub}^{2-\lambda_G}\), where
\(\lambda_G=12\tau_{\rm cub}G/(1+2\tau_{\rm cub}G)\).
The common power cancels in the completed first derivative; the actual
normalized section maps have ordinary \(C^3\) bounds. The general
raw cancellation is displayed in \eqref{sharp:upper-tail-ratio}.
Second, near an axis root the normal height has an action exponential.
The proof bounds derivatives on the product of the two separate section
domains, using this exponential to pay the losses in each root distance.
Only then is the estimate restricted to the actual matching graph.
This controls the two root distances independently and includes the
dependence on the moving entrance.

The other coefficient faces use the same physical chain rule with
different positive weights, as recorded in
Lemma~\ref{neg:lem-completed-curvature}. On the reversible face the
principal curvature vanishes. There the exact first integral permits
division by the coefficient distance before taking the limit, giving
the relative estimate \eqref{neg:relative-estimate}.
The coefficient cover selects one of these estimates for each whole
coefficient fibre.
